\section*{Synthese und Fazit}
\addcontentsline{toc}{section}{Synthese und Fazit}

StudyMummy koppelt die vier kognitiven Schichten in einem hierarchischen,
nachrichtengetriebenen Multi-Agent-System: Perception bildet den Lernzustand ab,
Sitzungen, Chatlogs und Vektoren tragen Memory, der Planner plant und der Tutor handelt.
Planner, Tutor und Reviewer besitzen getrennte Ziele, Capabilities und lokale
Zustände. Über typisierte Nachrichten führen Toolbeobachtungen und Reviewkritik
innerhalb eines Turns zu Revision oder Neuplanung. Der Regelkreis ist damit bis
zur konfigurierten Rundengrenze geschlossen.

Begrenzt bleibt das MAS durch serielle In-Process-Ausführung, einen gemeinsamen
LLM-Adapter und flüchtige Agentenzustände. Lernprofil und frühere Outcomes
steuern den Plan nicht. Schreibende Tools wirken vor dem Review und lassen sich
dort nicht zurückrollen. Vorrang haben deshalb ein transaktionaler Prüf- und
Commitpfad sowie danach ein Gedächtnis für validierte Outcomes. Die Tests
belegen Routing, Rückkopplung und Verträge, aber weder Tutorqualität noch
Lernwirksamkeit. Dafür fehlen Live-LLM- und Nutzerstudien.
